\Article{56}
Obwohl die Theorie quadratischer Reste und Nichtreste in der Arithmetik komplexer Zahlen in der Theorie biquadratischer Reste selbst enthalten ist, werden wir ihre Hauptsätze hier gesondert vortragen, bevor wir zu der letzteren übergehen. Der Kürze wegen behandeln wir jedoch nur den Hauptfall, in dem der Modul eine ungerade komplexe Primzahl ist.

Es sei $m$ ein solcher Modul und $p$ seine Norm. Offenbar kann eine ganze Zahl, die nicht durch $m$ teilbar ist, was hier stets vorauszusetzen ist, einem Quadrat nach dem Modul $m$ kongruent sein oder nicht, je nachdem, ob ihr Index bei einer als Basis genommenen primitiven Wurzel gerade oder ungerade ist. Im ersten Fall heißt diese ganze Zahl quadratischer Rest von $m$, im zweiten Nichtrest. Daraus folgt, dass unter $p-1$ Zahlen, die ein vollständiges System inkongruenter, durch $m$ nicht teilbarer Reste darstellen, die eine Hälfte zu den quadratischen Resten und die andere zu den quadratischen Nichtresten gehört. Jeder anderen Zahl außerhalb dieses Systems ist in dieser Hinsicht derselbe
\SourceLeaf{137}%
Charakter zuzuweisen wie der ihr kongruenten Zahl des Systems.

Ferner folgt daraus, dass das Produkt zweier quadratischer Reste ebenso wie das Produkt zweier Nichtreste ein quadratischer Rest ist, dagegen das Produkt eines quadratischen Restes und eines Nichtrestes ein Nichtrest. Allgemein ist ein Produkt aus beliebig vielen Faktoren ein quadratischer Rest oder Nichtrest, je nachdem, ob die Anzahl der Nichtreste unter den Faktoren gerade oder ungerade ist.

Zur Unterscheidung quadratischer Reste von Nichtresten bietet sich sogleich das folgende allgemeine Kriterium an:

Eine durch den Modul nicht teilbare Zahl $k$ ist ein quadratischer Rest oder Nichtrest desselben, je nachdem, ob $k^{\frac12(p-1)}\equiv1$ oder $k^{\frac12(p-1)}\equiv-1\pmod m$ gilt.

Die Richtigkeit dieses Satzes folgt unmittelbar daraus, dass bei einer beliebigen als Basis genommenen primitiven Wurzel der Index der Potenz $k^{\frac12(p-1)}$ entweder $\equiv0$ oder $\equiv\frac12(p-1)$ wird, je nachdem, ob der Index der Zahl $k$ gerade oder ungerade ist.

\Article{57}
Für einen gegebenen Modul lässt sich ein vollständiges System inkongruenter Reste leicht in zwei Klassen, quadratische Reste und Nichtreste, teilen; damit sind zugleich allen übrigen Zahlen ihre Klassen zugewiesen. Eine weit tiefere Untersuchung erfordert dagegen die Frage nach Kriterien, die jene Moduln, für die eine gegebene Zahl quadratischer Rest ist, von denen unterscheiden, für die sie Nichtrest ist.

Die reellen Einheiten $+1$ und $-1$ sind in der Arithmetik komplexer Zahlen tatsächlich Quadrate und daher auch quadratische Reste für \emph{jeden} Modul. Ebenso leicht folgt aus dem Kriterium des vorigen Artikels, dass $i$, und ebenso $-i$, quadratischer Rest jedes Moduls ist, dessen Norm $p$ die Form $8n+1$ hat, dagegen Nichtrest jedes Moduls, dessen Norm die Form $8n+5$ hat. Da es offenbar keinen Unterschied macht, ob $m$ oder eine der dazu assoziierten Zahlen $im,-m,-im$ als Modul genommen wird, dürfen wir voraussetzen, dass der Modul die primäre Zahl unter seinen Assoziierten ist (Artikel~36, II). Setzen wir also den Modul $=a+bi$, so dürfen wir $a$ als ungerade und $b$ als gerade annehmen. Dann gilt stets $aa\equiv1\pmod8$, während $bb$ entweder $\equiv0$ oder $\equiv4\pmod8$ ist, je nachdem, ob $b$ durch vier teilbar oder das Doppelte einer ungeraden Zahl ist. Somit sind $+i$ und $-i$ im ersten Fall offenbar quadratische Reste des Moduls, im zweiten Nichtreste.

\SourceLeaf{138}%
\Article{58}
Da die Entscheidung über den Charakter einer zusammengesetzten Zahl, ob sie quadratischer Rest oder Nichtrest ist, von den Charakteren ihrer Faktoren abhängt, genügt es offenbar, die Entwicklung von Kriterien zur Unterscheidung der Moduln, für die eine gegebene Zahl $k$ quadratischer Rest ist, von denen, für die sie Nichtrest ist, auf solche Werte von $k$ zu beschränken, die Primzahlen und überdies primär unter ihren Assoziierten sind. Bei dieser Untersuchung liefert das \emph{Untersuchen von Zahlenbeispielen} sogleich sehr elegante Sätze.

Beginnen wir mit der Zahl $1+i$, die sich als quadratischer Rest der Moduln
\[
\begin{gathered}
-1+2i,\ +3-2i,\ -5-2i,\ -1-6i,\ +5+4i,\ +5-4i,\\
-7,\ +7+2i,\ -5+6i,\quad\text{usw.}
\end{gathered}
\]
und als quadratischer Nichtrest der folgenden erweist:
\[
\begin{gathered}
-1-2i,\ -3,\ +3+2i,\ +1+4i,\ +1-4i,\ -5+2i,\\
-1+6i,\ +7-2i,\ -5-6i,\ -3+8i,\ -3-8i,\\
+5+8i,\ +5-8i,\ +9+4i,\ +9-4i\quad\text{usw.}
\end{gathered}
\]

Betrachten wir diese Übersicht aufmerksam, in die wir von je vier assoziierten Moduln stets den primären aufgenommen haben, so sehen wir leicht, dass alle Moduln $a+bi$ in der ersten Klasse $a+b$ $\equiv+1\pmod8$ erfüllen, die in der zweiten dagegen $a+b\equiv-3\pmod8$. Nehmen wir anstelle des primären Moduls $m$ den Assoziierten $-m$, so muss dieses Kriterium offenbar dahin geändert werden, dass für Moduln der ersten Klasse $a+b\equiv-1$ und für die der zweiten $\equiv+3\pmod8$ gilt. So ist allgemein, sofern uns die Untersuchung der Zahlenbeispiele nicht getäuscht hat und $a+bi$ eine Primzahl bezeichnet, in der $a$ ungerade. $b$ gerade ist, $1+i$ ein quadratischer Rest oder Nichtrest derselben, je nachdem, ob $a+b\equiv\pm1$ oder $\equiv\pm3\pmod8$ gilt.

Für die Zahl $-1-i$ gilt dieselbe Regel wie für $1+i$. Betrachten wir dagegen $1-i$ als Produkt aus $-i$ und $1+i$, so kommt $1-i$ offenbar derselbe Charakter zu, der $1+i$ beizulegen ist, wenn $b$ durch vier teilbar ist, dagegen der entgegengesetzte, wenn $b$ das Doppelte einer ungeraden Zahl ist. Daraus folgt leicht, dass $1-i$ ein quadratischer Rest der Primzahl $a+bi$ ist, wenn $a-b\equiv\pm1$ gilt, dagegen ein Nichtrest, wenn $a-b\equiv\pm3\pmod8$ gilt; dabei ist stets $a$ als ungerade und $b$ als gerade vorauszusetzen.

Dieser zweite Satz lässt sich auch aus dem ersten mithilfe eines allgemeineren Satzes ableiten, den wir folgendermaßen aussprechen:

\SourceLeaf{139}%
In der Theorie quadratischer Reste ist der Charakter der Zahl $\alpha+\beta i$ bezüglich des Moduls $a+bi$ derselbe wie der der Zahl $\alpha-\beta i$ bezüglich des Moduls $a-bi$.

Der Beweis dieses Satzes ergibt sich daraus, dass beide Moduln dieselbe Norm $p$ haben und dass, wenn $(\alpha+\beta i)^{\frac12(p-1)}-1$ durch $a+bi$ teilbar ist, auch $(\alpha-\beta i)^{\frac12(p-1)}-1$ durch $a-bi$ teilbar ist, während, wenn $(\alpha+\beta i)^{\frac12(p-1)}+1$ durch $a+bi$ teilbar ist, auch $(\alpha-\beta i)^{\frac12(p-1)}+1$ durch $a-bi$ teilbar sein muss.

\Article{59}
Gehen wir nun zu ungeraden Primzahlen über.

Wir finden, dass $-1+2i$ ein quadratischer Rest der Moduln $+3+2i$, $+1-4i$, $-5+2i$, $-5-2i$, $-1-6i$, $+7-2i$, $-3+8i$, $+5+8i$, $+5-8i$, $+9+4i$ usw. ist,

ein Nichtrest dagegen der Moduln $-1-2i$, $-3$, $+3-2i$, $+1+4i$, $-1+6i$, $+5+4i$, $+5-4i$, $-7$, $+7+2i$, $-5+6i$, $-5-6i$, $-3-8i$, $+9-4i$ usw.

Führen wir die Moduln der ersten Klasse auf ihre absolut kleinsten Reste nach dem Modul $-1+2i$ zurück, so finden wir nur $+1$ und $-1$, nämlich $+3+2i\equiv-1$, $+1-4i\equiv-1$, $-5+2i\equiv+1$, $-5-2i\equiv-1$ usw.

Dagegen erweisen sich alle Moduln der zweiten Klasse nach dem Modul $-1+2i$ entweder $+i$ oder $-i$ kongruent.

\Needspace{9\baselineskip}
Die Zahlen $+1,-1$ sind aber selbst quadratische Reste des Moduls $-1+2i$, während $+i$ und $-i$ dessen Nichtreste sind. Soweit man daher der Untersuchung von Zahlenbeispielen vertrauen darf, ergibt sich der folgende Satz: Die Zahl $-1+2i$ ist ein quadratischer Rest oder Nichtrest der Primzahl $a+bi$, je nachdem, ob diese ein quadratischer Rest oder Nichtrest von $-1+2i$ ist, sofern $a+bi$ unter den vier Assoziierten primär ist, oder vielmehr, wenn $a$ ungerade und $b$ gerade ist.

Aus diesem Satz folgen ohne Weiteres entsprechende Sätze über die Zahlen $+1-2i$, $-1-2i$, $+1+2i$.

\Article{60}
Untersuchen wir auf ähnliche Weise Zahlenbeispiele für $-3$ oder $+3$, so finden wir, dass beide quadratische Reste der Moduln $+3+2i$, $+3-2i$,
\SourceLeaf{140}%
$-1+6i$, $-1-6i$, $-7$, $-5+6i$, $-5-6i$, $-3+8i$, $-3-8i$, $+9+4i$, $+9-4i$ usw. sind,

Nichtreste dagegen von $-1+2i$, $-1-2i$, $+1+4i$, $+1-4i$, $-5+2i$, $-5-2i$, $+5+4i$, $+5-4i$, $+7+2i$, $+7-2i$, $+5+8i$, $+5-8i$ usw.

Die ersteren sind nach dem Modul $3$ einer der vier Zahlen $+1,-1,+i,-i$ kongruent, die letzteren dagegen einer von $+1+i,+1-i,-1+i,-1-i$. Jene sind gerade die quadratischen Reste der Zahl $3$, diese deren Nichtreste.

Diese Untersuchung von Zahlenbeispielen zeigt somit, dass die Primzahl \mbox{$a+bi$}, vorausgesetzt, dass $a$ ungerade und $b$ gerade ist, zur Zahl $-3$, und auch zu $+3$, dieselbe Beziehung hat wie diese zu jener, insofern eine ein quadratischer Rest oder Nichtrest der anderen ist.

Dehnen wir eine ähnliche Untersuchung auf andere Primzahlen aus, so finden wir dieses höchst elegante Reziprozitätsgesetz überall bestätigt und gelangen zu dem folgenden Fundamentalsatz über quadratische Reste in der Arithmetik komplexer Zahlen:

\emph{Bezeichnen $a+bi,A+Bi$ Primzahlen derart, dass $a,A$ ungerade und $b,B$ gerade sind, so ist entweder jede ein quadratischer Rest der anderen oder jede ein Nichtrest der anderen.}

Trotz der großen Einfachheit des Satzes ist sein Beweis jedoch mit großen Schwierigkeiten verbunden. Bei diesen halten wir uns hier nicht auf, da der Satz selbst nur ein besonderer Fall eines allgemeineren Satzes ist, der den Inhalt der Theorie biquadratischer Reste gleichsam ausschöpft. Zu dieser gehen wir daher nun über.

